% ============================================================
%  结语：因果的艺术与科学
%  1996 年 11 月 UCLA 教师研究讲座公开演讲
% ============================================================
\chapter*{结语：因果的艺术与科学}
\addcontentsline{toc}{chapter}{结语：因果的艺术与科学}
\markboth{结语：因果的艺术与科学}{结语：因果的艺术与科学}
\thispagestyle{plain}

\begin{center}
\kaishu\small
1996 年 11 月，作为加州大学洛杉矶分校（UCLA）\\
“教师研究讲座计划”（Faculty Research Lectureship Program）的一部分所作的公开演讲
\end{center}

\vspace{0.6em}

本讲的主题是因果性——即我们对世间何者导致何者的意识，以及这种意识为何要紧。

尽管因果性是人类思维中最基本的东西，它却是一个笼罩着神秘、争议与审慎的概念，因为科学家和哲学家一直难以界定：一个事件究竟在何时真正导致了另一个事件。我们都明白，公鸡打鸣并不会导致太阳升起；但即便是这样简单的事实，也不容易翻译成一个数学方程。

今天，我想与各位分享一组观念，我发现它们在研究这类现象时非常有用。这些观念已经催生出一些实用的工具，希望各位下次遇到因果问题时能用得上。

在座恐怕没有人不跟因果打交道。无论你是在评估双语教育项目的效果，还是在做小鼠如何辨别食物与危险的实验；无论你是在揣测尤利乌斯·凯撒当年为何渡过卢比孔河，还是在诊断病人，或是在预测谁将赢得总统选举——你面对的都是一张纠缠不清的因果之网。我下面要讲的故事，旨在帮助研究者应对这类考虑的复杂性，并厘清其含义。

本讲分为三部分。我先简略勾勒各个学科在因果问题上所遭遇困难的历史；然后概述一组观念，它们可以化解或消除其中若干历史难题；最后，出于对我工程学出身的敬意，我将说明这些观念如何导出简单而实用的工具，并在统计学与社会科学领域加以演示。

追根溯源，因果性最初似乎并不成问题。追问“为什么”的冲动，以及寻找因果解释的能力，在人类发展中很早就出现了。例如，《圣经》告诉我们，在尝了知识树上的果子之后不过几个小时，亚当就已经是因果论证的专家了。上帝问：“你是不是吃了那棵树上的果子？”亚当这样回答：“你所赐给我、与我同居的女人，她把那树上的果子给我，我就吃了。”夏娃同样娴熟：“那蛇引诱我，我就吃了。”

这个故事值得注意之处在于：上帝并没有要求解释，只要求事实——是亚当自己觉得有必要解释。信息很清楚：因果解释是人造的概念。

这个故事还有一点有趣之处：解释的用途，几乎仅限于推卸责任。的确，几千年来，解释没有别的功能。因此，能够使事情发生的，只有神、人和动物，而不是物件、事件或物理过程。

自然事件进入因果解释要晚得多，因为在古代世界里，事件本就是被注定的。风暴与地震由发怒的诸神掌管（幻灯片 2），它们本身无法为后果承担因果责任。即便是掷骰子（幻灯片 3）这样飘忽不定、不可预测的事件，也不被看作偶然，而被看作一道需要恰当解读的神谕。约拿先知就曾被这样一道神谕吓得魂飞魄散：他被指认为上帝面前的叛逃者，被抛入海中（幻灯片 4）。《约拿书》记载：“船上的人彼此说：‘来吧，我们掣签，看看这灾临到我们是因谁的缘故。’于是他们掣签，掣出约拿来。”

显然，在这艘豪华的腓尼基游轮上，“掣签”不是为了消遣，而是为了通信——一台处理生死攸关消息的单向调制解调器。

总之，在古代世界，因果力量的施动者要么是神祇——它们为着某种目的使事情发生；要么是人和动物——他们拥有自由意志，并因此受赏受罚。这种因果观念是幼稚的，但清晰而无碍。

问题照例始于工程：当人们必须造出机器来完成有用的工作之时（幻灯片 5）。随着工程师们野心渐长，他们认定地球也是可以撬动的（幻灯片 6），只不过不能靠一根杠杆。诸如此类规模的工程，需要由许多滑轮与轮子（幻灯片 7）组成的系统，一个驱动另一个。而一旦人们开始建造多级系统，因果性就发生了一件有趣的事：物理对象开始获得因果品格。

当这样一个系统出了故障，把责任推给上帝或操作者都是徒劳的——相反，一根断绳、一个生锈的滑轮才是更有用的解释，原因很简单：这些东西容易更换，换掉之后系统就能重新工作。从那一刻起，神和人不再是因果力量的唯一施动者——无生命的物件与过程也成了责任的分担者。一个轮子转动又停下，是因为前一个轮子转动又停下——操作者退居次要。

不出所料，这些新的因果施动者继承了其前辈——神与人——的某些特征。自然对象不仅承载功过，也承载力、意志，乃至目的。亚里士多德认为，只有依目的作出的解释，才是对“一物何以如此”的唯一完整而令人满意的解释。他甚至称之为终极因（final cause），即科学探究的最终目标。

自此，因果性承担了双重角色：一方面，原因是功过的归属对象；另一方面，原因又是物理控制流的载体。这种二元性在相对平静中存续下来（幻灯片 8），直到文艺复兴前后，才遭遇概念上的困难。

发生了什么，可以从雷科德（Recorde）《知识的城堡》（\textit{The Castle of Knowledge}）的扉页（幻灯片 9）上看出来。这是第一部英文科学著作，1575 年出版。命运之轮转动的力量，不是来自上帝的智慧，而是来自人的无知。随着上帝作为终极因的角色被人类知识接管，整个因果解释的观念便受到了攻击。

侵蚀始于伽利略（幻灯片 10）的工作。我们大多数人知道伽利略，是因为他因捍卫日心说而被押上宗教裁判所并遭监禁（幻灯片 11）。但就在这一切发生的同时，伽利略还悄然缔造了科学史上最深刻的一场革命。这场革命阐发于他 1638 年在远离罗马的莱顿出版的《关于两门新科学的对话》（\textit{Discorsi}，幻灯片 12），它由两条箴言构成：

其一，描述在先，解释在后——也就是说，“如何”先于“为何”；

其二，描述以数学的语言、即方程的语言来进行。

伽利略说，不要问一个物体下落，是因为它被下方吸引还是被上方推压。要问的是：你能多准确地预测物体走过某段距离所需的时间？这个时间又如何随物体不同、随斜面倾角改变而变化？而且，伽利略说，不要企图用人类语言中那些定性而滑溜的细微差别来回答这类问题；用数学方程的形式说出来（幻灯片 13）。

今天我们很难体会，这个想法在 1638 年听起来有多么古怪——那时距韦达（Vieta）引入代数记号还不到五十年。宣布代数是科学的通用语言，在今天听来，无异于宣布世界语是经济学的语言。大自然为什么偏偏同意说代数？在所有语言之中？可你没法跟成功争辩。物体走过的距离，果然正比于时间的平方。

比预测实验结果更成功的，是代数方程在计算方面的效用。它们使工程师有史以来第一次能够在“如果……会怎样”之外，提出“如何做”的问题。除了问“如果我们把梁收窄，它还能承重吗？”，他们开始问更难的问题：“要如何设计梁的形状，才能让它承重？”（幻灯片 14）这之所以可能，是因为有了求解方程的方法。代数机器不区分变量；我们不必用参数去预测行为，而可以把事情倒过来，用期望的行为去解出参数。

现在让我们集中谈伽利略的第一条箴言——“描述在先，解释在后”——因为科学家们把这一想法看得很重，它使科学的性质从思辨转向了经验。物理学中涌现出大量极其有用的经验定律。斯涅尔定律（幻灯片 15）、胡克定律、欧姆定律、焦耳定律，都是纯经验概括的例子：它们在得到更基本原理的解释之前很久，就已被发现并使用。

然而，哲学家们不愿放弃因果解释的观念，继续为那些成功的伽利略式方程寻找起源与依据。例如，笛卡儿把原因归于永恒真理；莱布尼茨把因果变成一条自明的逻辑法则。最后，在伽利略之后约一百年，一位名叫大卫·休谟的苏格兰哲学家（幻灯片 16）把伽利略的第一条箴言推到了极端（幻灯片 17）。

休谟令人信服地论证说，“为何”不仅次于“如何”，而且完全是多余的，因为它已被“如何”所吞并。在休谟《人性论》（幻灯片 18）第 156 页上，我们读到那段把因果性撼动得至今未能复原的文字。我每次读它都觉得过瘾：“因此我们记得曾经看见过我们称为火焰的那种对象，并且曾经感觉到我们称为热的那种感觉。我们同样回忆起它们在一切过去事例中的恒常结合。无需任何进一步的仪式，我们就把一个称为原因，把另一个称为结果，并由一个的存在推断另一个的存在。”

于是，在休谟看来，因果联系是观察的产物。因果观念是心灵的一种可习得习惯，其虚构程度近于视错觉，其短暂程度近于巴甫洛夫的条件反射。

很难相信休谟没有意识到他自己这套方案内在的困难。他很清楚，公鸡打鸣与日出处于恒常结合之中，却不导致日出；他也知道，气压计读数与降雨处于恒常结合之中，却不导致降雨。今天，这类困难被归入“伪相关”（spurious correlation）之名下，即“不蕴含因果的相关”。

现在，把休谟的箴言——一切知识都来自经验、并在心灵中被编码为相关——与我们关于“相关不蕴含因果”的观察合在一起，我们就得到了因果性的第一个谜题：人究竟如何获得因果知识？

公鸡的例子表明，前后相继的规则性并不充分；那么什么才是充分的？什么样的经验模式，才配得上把一种联系称为“因果的”？更进一步：什么样的经验模式，才会让人们相信一种联系是“因果的”？

如果说第一个谜题关乎因果联系的学习，第二个谜题则关乎它的使用：假如我告诉你某种联系是或不是因果的，这会造成什么差别？继续用我们的例子：假如我告诉你公鸡确实导致太阳升起，这会造成什么差别？

这听起来也许微不足道。显而易见的回答是：知道“什么导致什么”，会大大改变我们的行动方式。如果公鸡打鸣导致太阳升起，我们就可以让公鸡早些醒来、让它打鸣，从而把夜晚缩短——比如说，给它讲一个最新的公鸡笑话。

但这个谜题并不像看上去那么微不足道。如果因果信息除了前后相继的规则性之外还有经验含义，那么这种信息就应当出现在物理定律之中。可它并没有！

哲学家伯特兰·罗素 1913 年提出了这一论证（幻灯片 19）：“所有哲学家，”罗素说，“都以为因果性是科学的基本公理之一；然而奇怪的是，在发达的科学中，‘原因’一词从不出现……我相信，因果律是过时的遗物，像君主制一样残存至今，只是因为人们误以为它无害。”

另一位哲学家帕特里克·萨普斯（Patrick Suppes）曾为因果性的重要性辩护，他指出：“《物理评论》几乎每一期都至少有一篇文章，标题里用到‘原因’或‘因果性’。”

从这场交锋中我们得出的结论是：物理学家以一种方式谈话、写作、思考，却以另一种方式表述物理学。如果因果性仅仅被当作一种方便的通信装置——一种简写，用来表达那些否则需要许多方程才能写出的复杂物理关系模式——那么这种双语现象本可被原谅。说到底，科学中充满了缩写：我们说“把 $x$ 乘以 5”，而不说“把 $x$ 自身相加 5 次”；我们说“密度”，而不说“重量与体积之比”。为什么偏偏要挑因果性的毛病？

“因为因果性不一样，”罗素勋爵会这样反驳，“它不可能是缩写，因为物理定律全是对称的、双向成立的，而因果关系是单向的，从原因到结果。”

以牛顿定律为例：
\[
f = ma.
\]
代数的规则允许我们把这条定律写成五花八门的句法形式，而意义全都相同——只要知道三个量中的任意两个，第三个就确定了。然而在日常话语中，我们说力导致加速度，而不说加速度导致力，并且我们对这一区别感受强烈。同样，我们说比值 $f/a$ 有助于我们确定质量，而不说它导致质量。这些区别并不为物理方程所支持，这就引出一个问题：整套因果词汇是否纯属形而上学，只是“像君主制一样残存至今……”

幸运的是，极少有物理学家理会罗素的诘难。他们继续在办公室里写方程，在餐厅里谈因果；他们以惊人的成功轰开了原子，发明了晶体管和激光。工程学的情形也一样。但在另一个舞台上，这种张力无法被忽视，因为在那里，区分因果关系与其他关系的要求非常明确。这个舞台就是统计学。

故事始于大约一百年前相关（correlation）的发现。弗朗西斯·高尔顿（幻灯片 20）——指纹术的发明者、查尔斯·达尔文的表弟——很可以理解地想要证明才智与德行在家族中传承。高尔顿的研究促使他考虑各种方法，去度量一类个人或对象的属性与另一类属性之间的关联。1888 年，他测量了一个人前臂的长度和这个人脑袋的大小，并追问这两个量在多大程度上能相互预测（幻灯片 21）。他偶然发现了如下结果：如果把一个量对另一个量作图，并适当地缩放两个坐标轴，那么最佳拟合直线的斜率具有一些漂亮的数学性质。只有当两个量能够精确互测时，斜率才为 $1$；只要预测不比随机猜测更好，斜率就是 $0$；而最令人称奇的是，无论用 $X$ 对 $Y$ 作图还是用 $Y$ 对 $X$ 作图，斜率都相同。

“很容易看出，”高尔顿说，“这种共相关必定是两个器官的变异部分地由共同原因所致的结果。”

这里，我们第一次有了一个客观的度量，用来刻画两个变量如何“相关”；它严格基于数据，不掺人的判断或意见。

高尔顿的发现使他的一位弟子卡尔·皮尔逊（幻灯片 22）为之倾倒——此人如今被视为现代统计学的奠基人之一。皮尔逊当时 30 岁，已是一位有成就的物理学家和哲学家，即将转行做律师。45 年后（幻灯片 23），他这样描述自己最初对高尔顿发现的反应：

“我感觉自己像德雷克时代的海盗。……我把高尔顿的意思理解为：存在一个比因果更宽泛的范畴，即相关，而因果只是相关的极限；这种新的相关概念，把心理学、人类学、医学和社会学的大部分领域都带入了数学处理的范围。”

皮尔逊被人描述为“具有那种把汉尼拔带过阿尔卑斯山、把马可·波罗带到中国的干劲和决心”。当皮尔逊感觉自己像个海盗时，你可以确信他一定拿到了他的战利品。1911 年，他的《科学的语法》（\textit{The Grammar of Science}）第三版出版。书中新增了一章，题为“偶然性与相关——因果性的不充分”，皮尔逊在这一章里说：

“在‘物质’和‘力’这些已被弃置的基本概念之外，现代科学那不可索解的奥秘之中还矗立着另一个偶像，即原因与结果的范畴。”

那么，皮尔逊用什么来取代这个古老的原因—结果范畴呢？你简直不会相信自己的耳朵：列联表（contingency table，幻灯片 24）。

“这样的表称为列联表，而对两事物之间关系的描述，其最终的科学陈述总可以回溯到这样一张列联表。……一旦读者领会了这种表的性质，他也就把握了原因与结果之间关联这一概念的要旨。”

于是，皮尔逊断然否认在相关之外还需要一个独立的因果联系概念。他终生持此观点，因此在他的任何技术论文中都不提因果。他对“意志”“力”这类万物有灵论概念的攻击如此猛烈，对决定论的拒斥如此彻底，以致因果性还没来得及在统计学中扎根，就被他连根拔除了。

又过了 25 年，另一位意志坚强的人——罗纳德·费希尔爵士（幻灯片 25）——才让统计学家们表述出随机化实验（randomized experiment）：这是从数据检验因果关系的唯一得到科学证明的方法，而且至今仍是主流统计学所允许的唯一因果概念。

大体上，今天的局面依然如此。如果我们数一数关于因果性的博士论文、研究论文或教科书页数，得到的印象是：皮尔逊仍然统治着统计学。《统计科学百科全书》给“相关”分配了十二页，给“因果”只分配了两页——而这两页之中还有一页用来论证“相关不蕴含因果”。

让我们听听现代统计学家怎样谈论因果性。菲利普·道伊德（Philip Dawid），《生物计量学》（\textit{Biometrika}，皮尔逊创办的期刊）现任主编，承认：“因果推断是统计学全部问题中最重要、最微妙、也最被忽视的问题之一。”特里·斯皮德（Terry Speed），生物统计学会前主席（各位或许记得他是 O.\,J.~辛普森谋杀案审判中的专家证人），宣称：“对因果性的考虑，应当按统计学一贯的方式来对待：最好完全不考虑；若非考虑不可，则须极为谨慎。”戴维·考克斯爵士（Sir David Cox）与南妮·韦尔穆特（Nanny Wermuth）在几个月前出版的一本书中这样致歉：“我们在本书中没有使用‘因果的’或‘因果性’这些词……我们之所以谨慎，是因为很少能从单项研究中得出关于因果性的确凿结论。”

这种谨慎与回避的立场，使许多指望统计学提供指引的领域陷于瘫痪，尤其是经济学和社会科学。一位顶尖社会科学家在 1987 年说：“如果更多的研究者放弃思考和使用原因、结果这类术语，那将非常有益健康。”

这种局面会是一个人的功劳吗？即便是一个像皮尔逊那样的海盗？我怀疑。可是，统计学曾给世界贡献了假设检验与实验设计这样强有力的概念，我们又该如何解释它如此早就放弃因果性呢？

一个显而易见的解释当然是：因果比相关难测得多。相关可以在一次非受控研究中直接估计，而因果结论需要受控实验。但这过于简单了；统计学家并不容易被困难吓退，而孩子们在没有做过任何受控实验的情况下，也学会了因果联系。

我相信，答案在更深的地方，它与统计学的官方语言、即概率的语言有关。

这话也许会让在座一些人吃惊：概率论的词汇表里没有“原因”这个词；我们无法用概率的语言表达“泥泞不导致下雨”这句话——我们所能说的只是两者相互关联或相互依赖，意思是：若发现其一，便可预期其二。

自然，如果我们缺少一种语言来明确表达某个概念，就不能指望围绕这个概念发展出科学活动。科学发展要求知识能够从一个研究可靠地转移到另一个研究；而正如伽利略在 350 年前所示，这种转移要求一种形式语言的精确性与计算上的便利。

稍后我会回来讨论语言与记号的重要性，但首先我想用另一个领域的例子来结束这段历史巡礼——在那个领域里，因果性也吃了不少苦头。这次是计算机科学，即符号的科学。这个领域相对年轻，却极其强调语言与记号，因此或许能为这个问题提供一个有用的视角。

当研究者开始用计算机编码因果关系时，因果性的两个谜题被以新的活力唤醒。

设想你就是这个机器人（幻灯片 26），正试图弄明白厨房或实验室里发生了什么。在概念上，机器人面临的问题，与一个试图为国债建模的经济学家、或一个试图理解疾病传播的流行病学家所面临的问题是一样的。机器人、经济学家和流行病学家，都需要用有限的动作和带噪声的观测，从环境中追索出因果关系。这就把他们直接放到了休谟的第一个因果谜题面前：如何做到？

因果性的第二个谜题在机器人的世界里同样起作用。假设我们想走个捷径，把我们关于这个房间里因果所知的一切都教给机器人（幻灯片 27）。机器人应当如何组织并利用这些信息？

于是，因果性的两个哲学谜题，现在被翻译成了具体而实际的问题：机器人应当如何通过与环境的互动获得因果信息？机器人应当如何处理来自其创造者—程序员的因果信息？

同样，第二个谜题并不像看上去那么微不足道。罗素勋爵关于因果关系与物理方程不相容的警告，现在表现为一个明显的逻辑缺陷。例如，给定信息“如果草地湿了，那么下过雨”和“如果我们打碎这个瓶子，草地会湿”，计算机会推出“如果我们打碎这个瓶子，那么下过雨”（幻灯片 28）。这类程序错误浮现之迅速、之具体，使人工智能程序成为研究因果性细则的理想实验室。

这就把我们带到本讲的第二部分：如何通过把方程与图结合起来，解决因果性的第二个谜题；以及这一解决如何使第一个谜题不再那么可怕。

这一解决方案中有两个统摄性的观念：第一，把因果性视为干预（intervention）之下行为的总括；第二，用方程与图作为一种数学语言，因果思想可以在其中得到表示与操作。而要把两者结合起来，还需要第三个概念：把干预视为对方程施行的手术（surgery）。

让我们从一个广泛使用因果性、且从未在因果性上遇到麻烦的领域开始：工程学。

这里有一张电路图的工程图纸（幻灯片 29），它显示了电路中各信号之间的因果关系。电路由与门和或门组成，每一个都在输入与输出之间执行某种逻辑功能。让我们仔细考察这张图，因为它的简单与熟悉极具欺骗性。这张图实际上是科学最伟大的奇迹之一。它所能传达的信息，比数百万个代数方程、概率函数或逻辑表达式还要多。

\begin{center}
\fbox{\parbox{0.72\textwidth}{\centering\small 原书插图（从略）：由与门、或门构成的多级电路原理图，标有输入 $X$、$Y$ 与输出 $Z$ 之间的信号流向}}\end{center}

使这张图强大得多的，是它不仅能预测电路在正常条件下的行为，还能预测电路在数百万种异常条件下的行为。例如，给定这张电路图，我们可以轻易说出：若某个输入从 $0$ 变为 $1$，输出会是什么。这是正常的，很容易用一个简单的输入—输出方程表达。现在说异常的部分。我们同样可以说出：当我们把 $Y$ 置为 $0$、或把它接到 $X$、或把这个与门换成或门，或当我们执行这数百万种操作的任意组合时，输出会是什么。这张电路的设计者并没有预想到、甚至没有考虑过如此古怪的干预；然而奇迹般地，我们能够预测它们的后果。何以如此？这种表示能力从何而来？

它来自早期经济学家所称的自主性（autonomy）。也就是说，图中的门代表彼此独立的机制——改变其中一个而不改变另一个是容易的。这张图正是利用了这种独立性，用那些在干预之下保持不变的构件来描述电路的正常功能。

我在博尔特楼（Boelter Hall）的同事们一定在想：我为什么站在这里，把一件工程学上的琐事说得像世界第八大奇迹。我这样做有三个理由。第一，我想说明，在工程师视为理所当然的实践里，蕴含着大量未被发掘的智慧。第二，我想提醒经济学家和社会科学家注意这种图示方法的好处。他们时断时续地使用一种类似的方法已逾 75 年，称为结构方程建模与路径图；但近年来，他们听任代数上的便利压倒了图示表示，也压倒了它的好处。最后，在我看来，这些图捕捉到了因果性的精髓——预测异常事态与新操作的后果的能力。例如，在赖特（S.~Wright）的图中（幻灯片 30），我们可以预测：如果我们改变图中标为输入（$E$）的环境因素，甚至改变标为亲代与子代之间中间节点（$H$）的遗传因素，豚鼠一窝幼崽可能有什么样的毛色图案。这样的预测无法基于代数分析或相关分析作出。

\begin{center}
\fbox{\parbox{0.72\textwidth}{\centering\small 原书插图（从略）：赖特（S.~Wright, 1920）的豚鼠毛色路径图，环境输入 $E$ 经若干中间节点 $H$ 作用于子代毛色}}\end{center}

这样看待因果性，就解释了为什么科学家以如此的热情追求因果解释，也解释了为什么获得一个因果模型总伴随着一种“深刻理解”和“尽在掌控”的感觉。

深刻理解（幻灯片 31）意味着不仅知道事物昨天如何表现，还知道事物在新的假设情境下将如何表现，而控制只是这类情境之一。有趣的是，当我们拥有这种理解时，即便没有实际的控制手段，我们也会感到“尽在掌控”。例如，我们没有办法实际控制天体运动，但引力理论仍给我们一种理解与掌控之感，因为它提供了一张假设性控制的蓝图。我们可以预测未曾预料的新事件对潮汐的影响——比如月球被流星击中，或引力常数突然减半——而且同样重要的是，引力理论向我们保证：对人世间事物的寻常操作不会控制潮汐。难怪因果模型被视为试金石，用以区分深思熟虑的推理与反应式的、本能的回应。鸟类和猴子也许可以被训练去完成复杂任务，比如接好一根断线，但那需要试错训练。相反，深思熟虑的推理者无须尝试这些操作，就能预见新操作的后果。

让我们把电路图的一部分放大（幻灯片 32），以便理解为什么图能预测方程所不能预测的结果。我们也不妨从逻辑门换成线性方程（让在座各位都更自在些），并假定我们面对的系统只包含两个部件：一个乘法器和一个加法器。乘法器把输入乘以 $2$；加法器把输入加上 $1$。

\begin{center}
\begin{minipage}{0.30\textwidth}
\centering
$Y = 2X$\\[4pt]
$Z = Y + 1$
\end{minipage}%
\begin{minipage}{0.62\textwidth}
\centering
\begin{tikzpicture}[causal,node distance=11mm]
\node (x) {$X$};
\node[draw,minimum width=10mm,minimum height=6.5mm,right=of x] (m) {$\times 2$};
\node[right=of m] (y) {$Y$};
\node[draw,minimum width=10mm,minimum height=6.5mm,right=of y] (a) {$+1$};
\node[right=of a] (z) {$Z$};
\draw[->] (x) -- (m);
\draw[->] (m) -- (y);
\draw[->] (y) -- (a);
\draw[->] (a) -- (z);
\end{tikzpicture}
\end{minipage}
\end{center}

可是，这些方程与右边的图等价吗？显然不等！如果等价，那么我们把变量调换一下，所得的两个方程就应当与下面所示的电路等价：

\begin{center}
\begin{tikzpicture}[causal,node distance=11mm]
\node (z2) {$Z$};
\node[draw,minimum width=10mm,minimum height=6.5mm,right=of z2] (a2) {$-1$};
\node[right=of a2] (y2) {$Y$};
\node[draw,minimum width=10mm,minimum height=6.5mm,right=of y2] (m2) {$\div 2$};
\node[right=of m2] (x2) {$X$};
\draw[->] (z2) -- (a2);
\draw[->] (a2) -- (y2);
\draw[->] (y2) -- (m2);
\draw[->] (m2) -- (x2);
\end{tikzpicture}
\end{center}

但这两个电路是不同的。上面那个告诉我们：如果我们物理地操纵 $Y$，它会影响 $Z$；而下面那个表明：操纵 $Y$ 会影响 $X$，却对 $Z$ 毫无影响。再者，对我们的方程再做一些代数运算，可以得到另外两个方程，它们不指向任何结构；它们只是表示三个变量上的两个约束，却没有告诉我们这些变量如何相互影响。

让我们更仔细地考察我们判定物理操纵 $Y$ 之效果的那个心理过程——比如说，把 $Y$ 置为 $0$（幻灯片 33）。显然，当我们把 $Y$ 置为 $0$ 时，$X$ 与 $Y$ 之间的关系就不再由乘法器给出——现在有一个新机制控制 $Y$，$X$ 在其中没有发言权。在方程表示中，这相当于把方程 $Y = 2X$ 替换为新方程 $Y = 0$，并求解一组新的方程，得到 $Z = 1$。如果我们对表示下方模型的那对方程施行同样的手术，得到的解自然不同。第二个方程需要被替换，于是得到 $X = 0$，而 $Z$ 不受约束。

现在我们看到，这种干预模型如何导出一个关于因果性的形式定义：“若我们能通过操纵 $Y$ 改变 $Z$，也就是说，若在手术式地移除 $Y$ 的方程之后，$Z$ 的解依赖于我们为 $Y$ 代入的新值，则 $Y$ 是 $Z$ 的原因。”我们也看到，图在这个过程中何其关键。图告诉我们：当我们操纵 $Y$ 时，应当删除哪一个方程。而当我们把方程变换为代数等价的形式时（如屏幕下方所示），这一信息被彻底冲掉了。仅凭那对方程，不可能预测把 $Y$ 置为 $0$ 的结果，因为我们不知道该施行什么手术——根本不存在“$Y$ 的方程”这种东西。

总之，干预相当于对方程施行手术（由图引导），而因果关系意味着预测这种手术的后果。

这是一个超越物理系统的普遍主题。事实上，用“抹去”方程来建模干预的想法，最早由经济学家赫尔曼·沃尔德（Herman Wold）于 1960 年提出，但他的学说几乎从经济学文献中消失了。史书把这一神秘的消失归因于沃尔德的个性，但我倾向于认为原因更深：早期计量经济学家是非常谨慎的数学家；他们竭力保持代数的干净与形式化，不能同意让它被图之类的花招污染。而正如我们在屏幕上所见，没有图，手术操作在数学上毫无意义，因为它对我们书写方程的方式很敏感。

在阐述这一新数学运算的性质之前，让我先演示它对厘清统计学与经济学中的概念何其有用。

我们为什么偏爱受控实验而非非受控研究？假设我们想研究某种药物治疗对患有某种疾病的病人的康复效果。支配每个病人行为的机制，在结构上类似于我们之前看到的电路图。康复是治疗与其他因素（如社会经济状况、生活方式、饮食、年龄等）的函数。这里只画出了其中一个因素（幻灯片 34）。

\begin{center}
\begin{tikzpicture}[causal,node distance=16mm and 18mm]
\node[draw,rounded corners=2pt,inner sep=4pt,font=\small] (z) at (0,0) {背景因素 $Z$};
\node[draw,rounded corners=2pt,inner sep=4pt,font=\small] (x) at (-1.9,-1.5) {治疗 $X$};
\node[draw,rounded corners=2pt,inner sep=4pt,font=\small] (y) at (1.9,-1.5) {康复 $Y$};
\draw[->] (z) -- (x);
\draw[->] (z) -- (y);
\draw[->] (x) -- (y);
\end{tikzpicture}
\end{center}

在非受控条件下，治疗的选择由病人自己决定，且可能取决于病人的社会经济背景。这就造成一个问题：我们无法判断康复率的变化是归因于治疗，还是归因于那些背景因素。我们想做的是比较背景相同的病人，而这正是费希尔的随机化实验所实现的事情。如何实现？它实际上由两部分组成：随机化与干预。

干预意味着我们改变个体的自然行为：我们把受试者分成两组，称为治疗组与对照组，并说服受试者服从实验政策。我们把治疗分配给一些在正常情况下不会寻求治疗的病人，而给另一些本来会接受治疗的病人服用安慰剂。用我们的新词汇来说，这就是手术——我们切断一条功能连接，并用另一条替换它。费希尔的洞见在于：把新的连接接到一次随机掷币上，就能保证我们想切断的那条连接确实被切断。理由在于：人们假定，一枚随机的硬币不受任何我们能在宏观层面测量的东西影响——当然也包括病人的社会经济背景。

这幅图为普遍接受的随机化试验程序提供了一个有意义且形式化的依据。相形之下，我们的下一个例子则用手术观念指出一些广受接受的程序中的不足。

这个例子（幻灯片 35）涉及一位政府官员，他试图评估某项政策（比如税收）的经济后果。有意识地决定增税或减税，是对经济模型的一次手术，因为它改变了建立模型时所处的条件。经济模型是依据某段时期的数据建立的，而在这段时期里，税收曾随某些经济状况或政治压力而升降。然而，当我们评估一项政策时，我们希望比较相同经济状况下的备选政策——也就是说，我们希望切断那条在过去把政策与那些状况绑在一起的连接。在这种设定下，当然不可能把我们的政策接到掷币上、做一次受控实验；我们没有那个时间，而且实验尚未结束，经济可能已经被毁掉了。尽管如此，我们应当做的分析是：从受未切除模型支配的数据中，推断被切除模型的行为。

\begin{center}
\fbox{\parbox{0.72\textwidth}{\centering\small 原书插图（从略）：税收政策评估的经济模型，表示税率与经济状况之间的双向依赖}}\end{center}

我说“应当”，是因为你在任何经济学教科书里都找不到这样的分析。正如我先前提到的，沃尔德的手术观念在 1970 年代被逐出了经济学文献；而我能找到的一切关于政策分析的讨论，都假定被切除的模型贯穿始终。人们假定，税收在评估时处于政府控制之下，这就足以把税收自始至终当作外生变量；而事实上，税收在建模阶段是内生变量，只有在评估时才转为外生。当然，我并不主张恢复手术模型就能让政府一夜之间平衡预算，但这确实值得一试。

现在让我们考察，手术解释如何化解罗素的诘难——即因果关系的方向性与物理方程的对称性之间的冲突。物理方程确实是对称的；但当我们比较“$A$ 导致 $B$”与“$B$ 导致 $A$”这两个说法时，我们谈的并不是同一组方程。我们比较的是两个世界模型，它们由两组不同的方程表示：一组中，$A$ 的方程被手术式地移除；另一组中，$B$ 的方程被移除。罗素大概会在此打断我们：“事实上只有一个世界模型，即由全部物理方程合在一起给出的模型，你怎么能谈两个世界模型？”答案是：没错。如果你想把整个宇宙都纳入模型，因果性就消失了，因为干预消失了——操纵者与被操纵者失去了区别。然而，科学家极少把整个宇宙当作研究对象。在大多数情况下，科学家从宇宙中切出一块，并宣布这一块为“内”，即研究焦点；宇宙的其余部分则被视为“外”或背景，并用我们所谓的边界条件加以概括。这种对内与对外的选择，在我们看待事物的方式中造成了不对称；而正是这种不对称，使我们能够谈论“外部干预”，从而谈论因果性与因果的方向性。

这一点可以用笛卡儿那幅经典的图（幻灯片 36）来说明。作为一个整体，这个手—眼系统对因果一无所知。它不过是一团杂乱粒子与光子，竭力服从薛定谔方程，而该方程是对称的。然而，从中切出一块——比如物体那一部分（幻灯片 37）——我们就可以谈论手的运动导致这条光线改变角度。换一种切法，聚焦于大脑那一部分（幻灯片 38），瞧，现在反倒是光线导致手的运动——方向恰好相反。教训是：正是我们切割宇宙的方式，决定了我们赋予原因与结果的方向性。每一次科学探究都默含着这样的切割。在人工智能中，J.~McCarthy 称之为“限定”（circumscription）；在经济学中，限定相当于决定哪些变量被视为内生、哪些被视为外生——是在模型之内，还是在模型之外。

让我们总结方程模型与因果模型的本质区别（幻灯片 39）。两者都用一组对称的方程描述正常条件。但因果模型还包含三个额外的要素：(i) 内与外的区分；(ii) 一个假定：每个方程对应一个独立的机制，因而必须作为独立的数学语句加以保留；(iii) 干预，它被解释为对这些机制施行的手术。这使我们更接近于实现那个梦想：让因果性成为物理学中友善的一员。

但还缺一个要素：代数。我们前面讨论过，代数的计算便利对伽利略时代的科学家和工程师何其重要。我们能否指望这种代数便利也服务于因果性？让我换一种说法：我们所知的科学活动由两个基本部分组成：观察（幻灯片 40）与干预（幻灯片 41）。两者的结合就是我们所谓的实验室（幻灯片 42）——一个我们控制某些条件、观察另一些条件的地方。碰巧的是，标准代数很好地服务于观察部分，却迄今无益于干预部分。方程代数、布尔代数、概率演算无不如此——它们都为观察语句而非干预语句而设。

以概率论为例。如果我们想求出“在看见草地湿了的条件下下过雨”的概率，我们可以把问题写成一个形式语句：$P(\text{Rain}\mid\text{Wet})$，读作：给定 Wet，Rain 的概率（幻灯片 43）。竖线代表“给定我们看见”。我们不仅能把这个问句写成形式语句，还能用概率论的机器把该语句变换成其他表达式。在我们的例子里，只要觉得更方便或更有信息量，左边的语句就可以变换成右边的语句。

但假设我们问一个不同的问题：“如果我们把草地弄湿，下过雨的概率是多少？”我们甚至无法用概率的句法表达这个查询，因为竖线已经被占用，意为“给定我们看见”。我们可以发明一个新符号 \textit{do}，每当我们看到竖线后面跟着 \textit{do}，就读作“给定我们做”——但这无助于我们计算问题的答案，因为概率的规则不适用于这种新读法。我们凭直觉知道答案应当是什么：$P(\text{Rain})$，因为把草地弄湿不会改变下雨的概率。但这个直觉答案以及其他类似的答案，能否被机械地推导出来，以便在直觉失灵时安抚我们的思绪？

答案是肯定的，而这需要一门新的代数。首先，我们给新算子“给定我们做”指定一个符号。其次，我们找出操作含有这个新符号的语句的规则。我们所用的过程，类似于数学家当年发现标准代数规则的过程。

设想你是 16 世纪的一位数学家，如今已是加法代数方面的专家，而你迫切感到需要引入一个新算子——乘法，因为你厌倦了整天把一个数加到它自己身上（幻灯片 44）。你要做的第一件事，是给新算子指定一个符号：乘。然后你深入这个算子的含义，从中推出它的变换规则。例如，乘法的交换律可以这样推出，结合律也一样，等等。这些我们现在在中学就学到了。

以完全相同的方式，我们可以推出支配我们新符号 $\dox{\cdot}$ 的规则。我们有一门用于“看见”的代数，即概率论。我们有一个新算子，它焕然一新，含义由手术程序给了我们，非常清楚。演绎的大门已经打开，结果见下一张幻灯片（幻灯片 45）。请不要惊慌，我并不指望各位现在就读懂这些方程，但我想你们仍能感受到这门新演算的意味。它由三条规则组成，允许我们把涉及行动与观察的表达式变换为同类表达式：

\begin{itemize}[leftmargin=2.2em,itemsep=0.45em]
\item 规则 1（观察的删去与插入）：若 $(Y\indep Z\mid X,W)_{G_{\bar X}}$，则
\[
P(y\mid\dox{x},z,w)=P(y\mid\dox{x},w).
\]
\item 规则 2（行动与观察的互换）：若 $(Y\indep Z\mid X,W)_{G_{\bar X\underline{Z}}}$，则
\[
P(y\mid\dox{x},\dox{z},w)=P(y\mid\dox{x},z,w).
\]
\item 规则 3（行动的删去与插入）：若 $(Y\indep Z\mid X,W)_{G_{\bar X\overline{Z(W)}}}$，则
\[
P(y\mid\dox{x},\dox{z},w)=P(y\mid\dox{x},w).
\]
\end{itemize}

第一条允许我们忽略一个无关的观察，第三条允许我们忽略一个无关的行动；第二条允许我们用一个对同一事实的观察来替换一个行动。右边那些符号是什么？它们是图在变换合法时给我们的“绿灯”。下一个例子里我们就会看到它们发挥作用。

这就把我们带到本讲的第三部分，我将演示截至目前所讲的观念如何用来解决具有实际重要性的新问题。

考虑那个长达一个世纪的争论：吸烟对肺癌有何影响（幻灯片 46）。1964 年，卫生总监（Surgeon General）发布报告，把吸烟与死亡、癌症、尤其是肺癌联系起来。该报告基于非实验研究，其中发现吸烟与肺癌之间存在强相关；其主张是：所发现的相关是因果的——如果我们禁烟，癌症病例率将大致等于今天在非吸烟人群中观察到的水平。

这些研究遭到烟草业的猛烈攻击，其背后有一些非常著名的统计学家，其中包括罗纳德·费希尔爵士。他们的主张是：所观察到的相关也可以用一个模型来解释，其中吸烟与肺癌之间没有因果联系。相反，可能存在一种未观测的基因型，它同时导致癌症、并产生对尼古丁的先天天生渴求。用我们的记号，这一主张可形式地写作 $P(\text{Cancer}\mid\dox{\text{Smoke}})=P(\text{Cancer})$，意思是：让全体人口吸烟或戒烟，对癌症病例率没有影响。受控实验本可在两个模型之间作出裁决，但这种实验不可能进行（如今也不合法）。

\begin{center}
\begin{tikzpicture}[causal,node distance=14mm and 26mm]
% 面板 (a)
\node[draw,rounded corners=2pt,inner sep=4pt,font=\small] (a1) at (0,0) {吸烟};
\node[draw,rounded corners=2pt,inner sep=4pt,font=\small] (a2) at (4.0,0) {肺癌};
\draw[->] (a1) -- (a2);
\node at (2.0,-1.0) {\small (a)};
% 面板 (b)
\node[draw,dashed,rounded corners=2pt,inner sep=4pt,font=\small] (b0) at (6.6,0) {基因型};
\node[draw,rounded corners=2pt,inner sep=4pt,font=\small] (b1) at (6.6,-1.8) {吸烟};
\node[draw,rounded corners=2pt,inner sep=4pt,font=\small] (b2) at (9.4,-1.8) {肺癌};
\draw[->,dashed] (b0) -- (b1);
\draw[->,dashed] (b0) -- (b2);
\draw[->] (b1) -- (b2);
\node at (8.0,-2.8) {\small (b)};
\end{tikzpicture}
\end{center}

这些都是历史了。现在让我们进入一个假设的时代：双方代表决定会面，消除分歧。烟草业承认吸烟与癌症之间可能存在某种弱因果联系，健康界代表也承认可能存在与遗传因素的某些弱联系。于是他们画出这个组合模型，问题归结为：从数据中评估各条联系的强度。他们把这个查询提交给一位统计学家，答案立刻回来：不可能。意思是：无法从数据估计强度，因为任何数据都能完美地拟合这两个极端模型中的任何一个。于是他们放弃，决定照常继续政治斗争。分手之前，有人提议：也许我们可以测量一些辅助因素来解决分歧。例如，既然因果联系模型基于这样的理解——吸烟通过肺内焦油沉积的积累影响肺癌——那么我们也许可以测量样本个体肺内的焦油沉积量，这或许能提供量化各条联系所需的信息。双方都认为这是个合理的建议，于是向统计学家提交了一个新查询：假定有焦油沉积的中间测量，我们能否求出吸烟对癌症的效应？统计学家带回了好消息：可计算，而且解以闭式数学形式给出。如何做到？

\begin{center}
\begin{tikzpicture}[causal]
\node[draw,dashed,rounded corners=2pt,inner sep=4pt,font=\small] (g) at (0,1.7) {基因型};
\node[draw,rounded corners=2pt,inner sep=4pt,font=\small] (s) at (-2.2,0) {吸烟};
\node[draw,rounded corners=2pt,inner sep=4pt,font=\small] (t) at (0,0) {焦油沉积};
\node[draw,rounded corners=2pt,inner sep=4pt,font=\small] (c) at (2.2,0) {肺癌};
\draw[->,dashed] (g) -- (s);
\draw[->,dashed] (g) to[bend left=26] (c);
\draw[->] (s) -- (t);
\draw[->] (t) -- (c);
\end{tikzpicture}
\end{center}

统计学家接到问题，把它当作一个中学代数问题来处理：我们需要在假设性行动之下，从非实验数据——即从不含行动的表达式——计算 $P(\text{Cancer})$。或者说：我们需要从初始表达式中消去“\textit{do}”符号。消去的过程就像寻常解代数方程：在每一步（幻灯片 47），依据图的一个子图所许可，应用一条新规则，最终得到一个不含“\textit{do}”符号的公式，它表示一个可从非实验数据计算的表达式。

\begin{center}
\fbox{\parbox{0.72\textwidth}{\centering\small 原书插图（从略）：do 演算的逐步推导，每一步变换由图中相应子图（以高亮标出）授权}}\end{center}

你大概在想：这个推导是否解决了吸烟—癌症之争？答案是否定的。即便我们能拿到焦油沉积的数据，我们的模型也相当简陋，因为它基于某些双方可能都不认同的假定——例如，吸烟与肺癌之间不存在不经焦油沉积中介的直接联系。那样一来，模型就需要细化，我们最终可能得到一个含二十个以上变量的图。当有人对我们说“你没有把某某因素考虑进去”时，不必惊慌。恰恰相反，图欢迎这类新想法，因为往模型里添加因素和测量太容易了。现在已经有了简单的检验，使研究者只需瞥一眼图，就能判断我们能否计算一个变量对另一个变量的效应。

下一个例子说明，一个长期存在的问题如何仅凭图示手段解决——并由新代数加以证明。这个问题称为调整问题，或“协变量选择问题”，它代表辛普森悖论的实践一面（幻灯片 48）。

辛普森悖论最早由卡尔·皮尔逊于 1899 年注意到，它涉及一个令人不安的观察：两个变量之间的任何统计关系，都可能因在分析中纳入额外因素而反转。例如，你可能做一项研究，发现吸烟的学生成绩更高；然而，如果你对年龄作调整，那么在每个年龄组内，情况恰恰相反，即吸烟预示着更低的成绩。如果你进一步对父母收入作调整，你又会发现，在每个年龄—收入组内，吸烟预示着更高的成绩，诸如此类。

同样令人不安的是：从来没有人能告诉我们，哪些因素应当纳入分析。这样的因素现在可以通过简单的图示手段来识别。展示辛普森悖论的经典案例发生在 1975 年，当时加州大学伯克利分校因研究生招生中的性别偏见受到调查。在这项研究中，总体数据显示男性申请者的录取率更高；但按系分解后，数据显示出略微偏向录取女性申请者的倾向。解释很简单：女性申请者倾向于申请竞争更激烈的系，而在这些系里，男女的录取率都低。

为了说明这一点，设想一条渔船上有两张不同的网，一张大网眼的，一张小网眼的（幻灯片 49）。一群鱼游向渔船，试图穿过。雌鱼去挑战小网眼，雄鱼则走容易的路。雄鱼穿过去了，只有雌鱼被抓住。从最终的渔获看，对雌鱼的偏好显而易见。然而，如果分开分析，每一张网无疑都更容易捕获雄鱼而非雌鱼。

另一个例子涉及一场称为“反向回归”（reverse regression）的争论，它在 1970 年代占据了社会科学文献。在薪资歧视案件中，我们应当比较同等资历的男女的薪资，还是应当比较同等薪资的男女的资历？值得注意的是，两种选择导致了相反的结论。结果发现，男性比同等资历的女性挣得更多；与此同时，男性又比同等薪资的女性资历更高。其寓意是：一切结论都对我们选择在比较时固定哪些变量极为敏感；这就是为什么调整问题在观察研究的分析中如此关键。

考虑一项观察研究，我们想求出 $X$ 对 $Y$ 的效应，例如治疗对反应（幻灯片 50）。我们可以想到许多与此问题相关的因素；有些受治疗影响，有些影响治疗，还有些同时影响治疗与反应。其中一些因素可能无法测量，如遗传特征或生活方式；另一些则可测量，如性别、年龄和薪资水平。我们的问题是：从这些因素中选出一个子集加以测量和调整，使得在相同测量值之下比较受试者并取平均时，我们得到正确的结果。

\begin{center}
\fbox{\parbox{0.72\textwidth}{\centering\small 原书插图（从略）：观察研究中 $X$ 对 $Y$ 之效应的调整问题，图中画出受处理影响、影响处理以及同时影响两者的各类因素}}\end{center}

让我们一起来走一遍检验两个候选测量 $Z_1$、$Z_2$ 是否充分的步骤（幻灯片 51）。步骤相当简单，即便在大图上也可以手工完成。不过，为了让各位感受其可机械化程度，我会讲得比较快。开始吧（幻灯片 52—56）。

\begin{center}
\fbox{\parbox{0.72\textwidth}{\centering\small 原书插图（从略）：调整准则的逐步检验，共六幅，逐次删去图上的节点与边，直至 $X$ 与 $Y$ 被分离}}\end{center}

在这些操作的最后，我们得到了问题的答案：“若 $X$ 与 $Y$ 被分离，则 $Z_1$ 与 $Z_2$ 是合适的测量。”

现在我简要总结本讲的核心信息。检验原因与结果确实困难；发现结果的成因更困难。但因果性并不神秘，也不是形而上学。它可以用简单的过程来理解，也可以用一种友善的数学语言来表达，随时可供计算机分析。

今天我向各位呈现的，是一种袖珍计算器、一把算盘（幻灯片 57），帮助我们以数学的精确性研究某些因果问题。这并没有解决因果性的全部问题，但符号与数学的力量不应被低估（幻灯片 58）。几个世纪以来，许多科学发现之所以被延误，就是因为缺少一种能够放大思想、并让科学家交流结果的数学语言。我确信，在我们这个世纪，许多发现之所以被延误，是因为缺少一种能够处理因果性的数学语言。例如，我确信，卡尔·皮尔逊如果在 1901 年就允许因果图进入他的数学，他本可以想出随机化实验这个主意。

但真正具有挑战性的问题还在前头：我们仍然没有对贫困、癌症和不宽容的因果理解，只有数据的积累与伟大头脑的洞见，最终才会带来这样的理解。数据遍地都是，洞见属于你们，而现在，一把算盘也交到了你们手上。我希望这一组合能放大其中每一个成分。

谢谢。

\section*{致谢}

幻灯片 1（丢勒，《亚当与夏娃》，1504 年铜版画）承蒙哈佛大学艺术博物馆福格艺术博物馆惠允，系威廉·格雷捐赠、出自弗朗西斯·考利·格雷收藏。摄影：Rick Stafford；图像版权 © President and Fellows of Harvard College, Harvard University。幻灯片 2（多雷，《罗得逃离》）版权归 William H. Wise \& Co.。幻灯片 3（埃及壁画：涅斐龙佩玩棋盘游戏）承蒙芝加哥大学东方研究所惠允。

下列图像复制自古籍书商目录，承蒙伦敦 Bernard Quaritch Ltd.\ 惠允：幻灯片 4、5、6、7、8、9、15、27、31、36、37、38、40、42 及 58。

幻灯片 10 与 11 版权归 The Courier Press；幻灯片 13 与 14 经 Macmillan Library Reference USA 许可重印，出自 I.~Bernard Cohen 的 \textit{The Album of Science}，版权 © 1980 Charles Scribner's Sons。

幻灯片 16 承蒙加州州立大学长滩分校图书馆惠允。幻灯片 20 与 22 经剑桥大学出版社许可重印。幻灯片 25：摄影版权归 A.~C.~Barrington Brown，经许可复制。

幻灯片 30：取自 S.~Wright (1920)，载 \textit{Proceedings of the National Academy of Sciences} 第 6 卷；经美国哲学学会与芝加哥大学出版社许可复制。幻灯片 57 经 Vandenhoeck \& Ruprecht 与 The MIT Press 许可重印。

注：幻灯片 19、26、28—29、32—35 及 43—56 的彩色版本可从下列地址下载：
\begin{center}
\url{http://www.cs.ucla.edu/~judea/}
\end{center}
